% Generated draft: requires mathematical and source review.
\documentclass[11pt]{article}
\usepackage[margin=25mm]{geometry}
\usepackage{fontspec}
\setmainfont{Latin Modern Roman}
\usepackage{amsmath,amssymb,mathtools,unicode-math}
\setmathfont{Latin Modern Math}
\usepackage{enumitem,needspace,xurl,hyperref}
\setlength{\parindent}{0pt}
\setlength{\parskip}{5pt}
\setlength{\emergencystretch}{3em}
\newcommand{\N}{\mathbb{N}}
\newcommand{\Z}{\mathbb{Z}}
\newcommand{\Q}{\mathbb{Q}}
\newcommand{\R}{\mathbb{R}}
\newcommand{\C}{\mathbb{C}}
\newcommand{\F}{\mathbb{F}}
\newcommand{\A}{\mathbb{A}}
\newcommand{\PP}{\mathbb{P}}
\newcommand{\E}{\mathbb{E}}
\newcommand{\eps}{\varepsilon}
\newcommand{\dd}{\,\mathrm{d}}
\begin{document}
\section*{UnsolvedMath ID: 30005891}
\textbf{Source problem code:} catalogue entry
\subsection{Definitions and mathematical statement}
Let \(G\) be a finite group. A subset \(A\subseteq G\) is \emph{product-free} if
\[
\forall a,b\in A,\qquad ab\notin A.
\]
Define its density to be \(|A|/|G|\), and write
\[
\alpha_{\mathrm{pf}}(G)=\max\left\{\frac{|A|}{|G|}:A\subseteq G\text{ is product-free}\right\}.
\]
For an action of \(G\) on a finite set \(X\), a point \(x\in X\), and a subset \(I\subseteq X\), define
\[
K_{x,I}=\{g\in G:g\cdot x\in I\text{ and }g\cdot I\subseteq X\setminus I\}.
\]
The source presents these sets as examples of product-free sets. It imposes no explicit faithfulness or transitivity condition on the action and no additional restriction on \(x\) or \(I\).

The precise question attributed to N. Lifshitz is Question 10 of the supplied excerpt. Its representation-theoretic hypothesis requires clarification: literally, every finite group has a one-dimensional trivial irreducible representation, so a minimum over all irreducible representations cannot exceed one. Under the conventional interpretation that the minimum is taken over nontrivial irreducible complex representations, put
\[
D(G)=\min\{\dim_{\C}V:V\text{ is a nontrivial irreducible complex representation of }G\}.
\]
The exclusion of the trivial representation and the choice of \(\C\) are interpretive qualifications, not explicit hypotheses in the supplied excerpt.

Under that interpretation, the intended question is whether there exist absolute constants \(C,c>0\) such that, for every finite group \(G\) for which \(D(G)\) is defined and \(D(G)>1\), there exist a finite \(G\)-set \(X\), a point \(x\in X\), and a subset \(I\subseteq X\) satisfying
\[
|X|\leq C D(G)^c,
\qquad
\frac{|K_{x,I}|}{|G|}=\alpha_{\mathrm{pf}}(G).
\]
Here the bound expresses the source notation \(D^{O(1)}\); its intended uniformity should be checked against the report. The assertion concerns the existence of a maximum-cardinality product-free set of this form, rather than an explicit classification of all maximizers. The registry summary additionally asks to determine maximum density and structure in finite quasirandom groups, which is broader than the displayed existence question.

The excerpt contains a detached continuation involving an undefined function \(f(n)\). It attributes to reference [7, Corollary 15] the existence of an unspecified set for \(f(n)<\exp((\log n)^k)\), for any fixed \(k\); asserts nonexistence for \(f(n)\gg n!\), citing an unbounded minimal number of generators; and suggests the possibility \(f(n)\approx\exp(n^{0.001})\). Neither \(f\), \(n\), the set, nor their connection with Question 10 is supplied. These fragments cannot be incorporated as hypotheses or established consequences of the question above.
\subsection{Short English statement}
For a finite group whose nontrivial irreducible complex representations all have large dimension, the intended question asks whether some largest product-free subset comes from a group action on a set of size polynomial in the minimum such dimension. The representation convention and the separation of Question 10 from adjacent text require review. The supplied registry labels the problem partially solved; that assessment is not independently verified here.
\subsection{Sources}
\begin{itemize}
\item Supplied local UnsolvedMath registry record, permanent ID \(30005891\), title \emph{Extremal Product-Free Sets in Quasirandom Groups}, exact source code \texttt{catalogue entry}. The record attributes Question 10 to N. Lifshitz and gives the proposed year as 2024; its structured proposer field is empty.
\item \emph{Mini-Workshop: Growth and Expansion in Groups} (2024), Oberwolfach Reports, pp.~1005--1040, DOI 10.4171/owr/2024/17: \url{https://doi.org/10.4171/owr/2024/17}. This report citation is supplied provenance and has not been independently inspected here.
\item Supplied literature link, labelled ``largest product-free subsets'': \url{https://doi.org/10.1007/s00222-024-01273-1}.
\item Supplied literature link, labelled ``product mixing'': \url{https://arxiv.org/abs/2401.15456}.
\item Supplied literature link, labelled ``OWR report'': \url{https://ems.press/content/serial-article-files/49475}.
\item Literature assessment dated exactly \texttt{2026-08-22}: the supplied curation labels the problem partially solved and states that the proposed form is confirmed for alternating groups of sufficiently large degree, with Kedlaya-type extremal families, while the universal assertion for finite groups with minimum representation dimension \(D\) remains open. This is supplied dated curation derived from the OWR record, not an independently verified literature assessment; its status and cited claims require independent checking.
\item The excerpt attributes the definition of product-free sets to Babai and S\'os (1985), reference [4], and cites [7, Corollary 15] in the detached continuation. Full bibliographic details for these numbered references are not supplied.
\end{itemize}
\end{document}
